\chapter{推理时计算缩放：用更多思考换更好答案}
\label{ch:ttc}
%% ============================================================

传统的 LLM 使用方式是：
输入一个问题，模型立刻给出答案。
但 OpenAI 的 o1 系列和 DeepSeek-R1 开创了一种新范式：
\textbf{让模型在回答之前先「思考」}，
生成大量的推理过程（thinking tokens），
然后才给出最终答案。

\section{核心洞察}

\begin{corebox}[推理时计算的颠覆性影响]
作为一个粗略的估计：
\textbf{7B 模型 + 100$\times$ 推理计算 $\approx$ 70B 模型标准推理}
（基于 Snell et al.\ (2024)~\cite{snell2024scaling} 的研究）。
这意味着你可以用一个便宜的小模型，
通过投入更多推理计算来获得与大模型相当的质量。
这彻底改变了「模型大小 = 能力」的等式。
\end{corebox}

传统的范式是\textbf{训练时缩放}（train-time scaling）：
投入更多计算来训练更大的模型，
推理时每个问题使用固定的计算量。
Scaling Laws 描述的就是这种范式。

推理时计算缩放代表了一种\textbf{互补的范式}：
模型大小固定，
但根据问题的难度动态调整推理时的计算量。
简单问题快速回答（几个 token 的思考），
复杂问题深度推理（数千个 token 的思考）。

这两种范式的关系可以类比为人类学习和思考的区别：
训练时缩放相当于\textbf{学更多年}（更多教育），
推理时缩放相当于\textbf{想更久}（更深思考）。
一个博士可能比一个本科生知道更多（训练时缩放），
但本科生如果花足够多的时间仔细思考，
在某些具体问题上也可能得到更好的答案（推理时缩放）。

但这个类比也揭示了推理时缩放的\textbf{边界}。
“想更久”只在你\textbf{有相关知识}时才有用，
一个不懂量子力学的学生，
无论思考多久也解不出量子力学的习题。
这个直觉在 LLM 推理中同样适用，
并且可以更精确地刻画。

\textbf{推理时缩放有效的任务类型}：
\begin{itemize}[noitemsep]
  \item \textbf{多步推理}（数学证明、逻辑推导）：
        更多的 thinking tokens 允许模型展开更长的推理链，
        检查更多可能的推导路径。收益最大。
  \item \textbf{代码生成与调试}：
        模型可以“在脑中运行”代码来检查逻辑，
        更多的推理时间 = 更多的心智模拟步骤。
  \item \textbf{复杂规划任务}：
        Agent 需要在多步行动中权衡利弊，
        更多的推理使得规划更周全。
  \item \textbf{需要考虑多种情况的决策}：
        法律分析（考虑多个先例）、
        医疗诊断（排除多种可能）。
\end{itemize}

\textbf{推理时缩放效果有限的任务类型}：
\begin{itemize}[noitemsep]
  \item \textbf{纯事实回忆}：
        “法国的首都是什么？”：
        模型要么知道，要么不知道，
        更多的思考不会“发明”新知识。
  \item \textbf{感知类任务}：
        图像描述、语音识别，
        这些任务的信息在输入中，
        而非需要通过推理发现。
  \item \textbf{简单指令遵循}：
        “把这段话翻译成法语”，
        任务结构清晰，标准推理已经足够。
  \item \textbf{极度开放的创意任务}：
        “写一首关于春天的诗”，
        没有“正确答案”，
        更多的推理可能反而导致“过度思考”
        而失去创意的自然流动。
\end{itemize}

这种“有效/无效”的划分有一个统一的解释：
推理时缩放在\textbf{搜索空间大、但答案可验证}的任务中最有效。
大的搜索空间意味着需要探索多条路径（推理时计算的价值所在），
可验证意味着能区分好的和坏的路径（避免推理变成随机游走）。
数学和代码完美满足这两个条件，
搜索空间指数级大，但答案可以精确验证。
纯事实回忆搜索空间为零（答案唯一且已知），
开放创意的答案不可验证，
两者都不适合推理时缩放。

\begin{whybox}[为什么推理时缩放可能比训练时缩放更高效？]
Snell et al.\ (2024) 的研究给出了一个关键发现：
对于许多具体问题，
将计算预算花在推理阶段比花在扩大模型更具\textbf{计算效率}。

直觉解释：
训练时扩大模型，所有问题都受益于更多的参数，
但简单问题本来就能做对，额外的参数是浪费。
推理时缩放允许\textbf{按需分配}，
简单问题快速回答节省计算，
难的问题多花时间保证质量。
这种动态分配本质上是一种更优的计算资源配置。

类比：
训练时缩放像「给所有学生都上博士课程」，
对已经懂的学生是浪费；
推理时缩放像「允许每个学生根据题目难度决定思考时间」，
更灵活，更高效。
\end{whybox}


\section{推理时缩放的数学框架}

\subsection{Snell et al.\ 的计算最优缩放理论}

Snell et al.~\cite{snell2024scaling} 的核心贡献是建立了
推理时计算的\textbf{缩放法则}（scaling law），
回答了一个关键问题：
给定固定的推理计算预算 $C$，
如何\textbf{最优分配}来最大化性能？

他们分析了两种缩放机制：

\textbf{机制一：基于验证器的搜索}。
生成 $N$ 个候选回答，
用 Process Reward Model（PRM）评分，
选择最优。
性能提升可以建模为：
\begin{equation}
  P(\text{correct} \mid N, \text{PRM}) = f(N, q_{\text{PRM}})
\end{equation}
其中 $q_{\text{PRM}}$ 是验证器的质量，
$N$ 是候选数量。
关键发现是：\textbf{PRM 质量比候选数量更重要}。
一个好的 PRM + 少量候选（$N = 4$）
通常优于一个差的 PRM + 大量候选（$N = 256$）。

\textbf{机制二：自适应分布更新}。
根据具体问题动态调整模型的回答策略。
对于推理模型，这体现为
根据问题难度生成不同长度的推理链。

\textbf{最优分配}的核心结论是：
\begin{equation}
  \text{Optimal strategy}(C, d) = \begin{cases}
    \text{少量深度推理} & \text{如果 } d \text{ 高（难题）} \\
    \text{多次浅层尝试} & \text{如果 } d \text{ 低（简单题）} \\
    \text{放弃推理时缩放} & \text{如果 } d \text{ 极低或极高}
  \end{cases}
\end{equation}
其中 $d$ 是问题的难度。
这个结论有重要的实践意义：
\begin{itemize}[noitemsep]
  \item 对于「模型本来就能做对」的简单问题，
        额外的推理计算几乎没有收益，直接回答即可。
  \item 对于「中等难度」的问题，推理时缩放的收益最大，
        相当于用计算换取了更大模型的效果。
  \item 对于「远超模型能力」的极难问题，
        再多的推理计算也无济于事，
        模型不具备必要的知识或推理能力。
\end{itemize}

\begin{whybox}[推理时缩放的数学等价性]
Snell et al.\ 的一个惊人发现是\textbf{计算等价性}：
一个较小的模型（如 Llama 3.1 8B）
通过最优的推理时计算分配，
可以在某些问题上匹配甚至超过
一个大 14$\times$ 的模型（Llama 3.1 70B）的性能。

形式化地说，对于难度为 $d$ 的问题集合，
存在一个推理计算预算 $C^*(d)$ 使得：
\begin{equation}
  \text{Acc}(\text{小模型}, C^*(d)) \geq \text{Acc}(\text{大模型}, C_0)
\end{equation}
其中 $C_0$ 是标准推理的基线计算量。
当 $C^*(d) < C_{\text{train}}$（训练大模型的额外成本摊到推理上）时，
推理时缩放就是更经济的选择。

但这个等价性不是\textbf{普适}的：
它高度依赖问题难度。
对于极难的问题，没有任何推理计算量
能让小模型匹配大模型。
这就是为什么推理时缩放是训练时缩放的\textbf{互补}而非\textbf{替代}。
\end{whybox}

\subsection{三种 Scaling Laws 的统一视角}

到 2026 年，LLM 领域已经形成了
\textbf{三种 scaling laws} 的统一框架：

\begin{enumerate}[noitemsep]
  \item \textbf{预训练缩放}（Kaplan 2020, Chinchilla 2022）：
        更多的参数 $\times$ 更多的数据 $\times$ 更多的训练计算
        → 更好的基础能力。
        描述的是\textbf{模型知识}的缩放。
  \item \textbf{后训练缩放}（RL、RLHF、SFT）：
        更多的对齐训练 → 更好地遵循指令和推理。
        描述的是\textbf{模型行为}的缩放。
  \item \textbf{推理时缩放}（Snell 2024, o1/R1）：
        更多的推理计算 → 在具体问题上更好的表现。
        描述的是\textbf{问题求解}的缩放。
\end{enumerate}

这三种缩放在\textbf{效率前沿}上各有优势：
预训练缩放适合需要广泛知识的场景，
后训练缩放适合需要特定行为模式的场景，
推理时缩放适合需要深度推理的特定问题。
2026 年的趋势是同时优化这三个维度：
所谓的\textbf{效率缩放}（efficiency scaling），
目标是用 \$1M 的总成本达到以前需要 \$100M 的效果。

\begin{figure}[ht]
\centering
\begin{tikzpicture}[
  >={Stealth[length=2.5pt,width=2.5pt]},
  every node/.style={font=\small},
]
% Axes
\draw[->, thick] (0, 0) -- (10.5, 0) node[right] {推理计算量 ($\times$ 基线)};
\draw[->, thick] (0, 0) -- (0, 7.2) node[above] {准确率 (\%)};

% Axis labels (log2 scale: 2 units per doubling)
\foreach \x/\lab in {0/1, 2/2, 4/4, 6/8, 8/16, 10/32} {
  \draw (\x, -0.1) -- (\x, 0.1) node[below=2pt] {$\lab\times$};
}
\foreach \y/\lab in {1/30, 2/45, 3/60, 4/75, 5/85, 6/95} {
  \draw (-0.1, \y) -- (0.1, \y) node[left=2pt] {\lab};
}

% Easy problems curve (saturates fast)
\draw[very thick, figGreen] plot[smooth, tension=0.7] coordinates
  {(0, 5.0) (2, 5.5) (4, 5.8) (6, 5.92) (8, 5.96) (10, 5.98)};
\node[figGreen, right] at (10.1, 5.98) {简单题};

% Medium problems curve (sweet spot at 4x--16x)
\draw[very thick, figBlue] plot[smooth, tension=0.7] coordinates
  {(0, 2.2) (2, 3.4) (4, 4.5) (6, 5.1) (8, 5.4) (10, 5.55)};
\node[figBlue, right] at (10.1, 5.55) {中等题};

% Hard problems curve (diminishing returns)
\draw[very thick, figRed] plot[smooth, tension=0.7] coordinates
  {(0, 0.5) (2, 1.0) (4, 1.7) (6, 2.2) (8, 2.5) (10, 2.8)};
\node[figRed, right] at (10.1, 2.8) {极难题};

% Annotation: sweet spot (4x--16x range)
\draw[dashed, figGray] (4, 0) -- (4, 6.3);
\draw[dashed, figGray] (8, 0) -- (8, 6.3);
\node[figGray, font=\footnotesize, align=center, fill=white, inner sep=2pt] at (6, 6.8)
  {“甜蜜区间”：\\中等题获益最大};

% Annotation: diminishing returns
\draw[->, thick, figOrange] (8.5, 1.5) -- (9.5, 1.5);
\node[figOrange, font=\footnotesize, left] at (8.5, 1.5) {收益递减};
\end{tikzpicture}
\caption{推理时计算缩放曲线示意。
简单题（绿色）很快饱和，额外计算浪费；
中等难度（蓝色）在 4--16$\times$ 计算量区间获益最大；
极难题（红色）收益始终有限，模型缺乏必要知识。
2--8$\times$ 基线计算量是性价比最优的“甜蜜区间”。}
\label{fig:ttc-scaling-curve}
\end{figure}

\textbf{Budget forcing}（预算强制）
是推理时缩放的一个实用技术：
根据问题的难度\textbf{预先分配}不同的计算预算。
简单问题限制 thinking tokens 在 256 以内，
中等问题分配 2K--8K，
困难问题分配 16K--64K。
这避免了模型在简单问题上「过度思考」浪费资源，
也保证了困难问题有足够的推理深度。


\section{三种策略}

推理时计算缩放有三种基本策略，
各自适用于不同的场景。

\subsection{策略一：扩展推理链（Sequential Scaling）}

\textbf{顺序缩放}让模型在内部进行长链推理，
分解问题、逐步推导、自我验证、纠正错误。
DeepSeek-R1~\cite{deepseekr1} 的 thinking tokens 就是这种策略。
每个 thinking token 都是在「推理计算」上的投资。
模型可能生成数千个 thinking token，
但最终答案的质量远高于直接回答。

一个典型的推理过程可能是这样的：
模型先分析问题结构，
然后尝试一种解法，
发现行不通后自我纠正（“wait, that's not right...”），
换一种方法，
验证中间结果，
最终得到答案。
这种\textbf{自我对话和纠错能力}
是 RL 训练的产物，不是人类预先编程的。

\textbf{思维链推理}（Chain-of-Thought, CoT）
是顺序缩放的基础形式。
即使不是专门训练的推理模型，
简单地在 prompt 中加上
“Let's think step by step” 就能提升推理质量。
但专门训练的推理模型（如 R1、o1）
的推理深度和自我纠错能力远超简单的 CoT prompting。

\textbf{训练推理模型}的核心方法是
GRPO~\cite{shao2024grpo}（Group Relative Policy Optimization）
配合\textbf{结果导向奖励}（outcome-based reward）。
训练过程不告诉模型「如何推理」，
只告诉它「答案对不对」，
模型必须自己发现有效的推理策略来最大化奖励。
这种纯粹由 RL 驱动的训练方式解释了为什么
R1 会自发地「发明」推理策略：
它不是在模仿人类的推理范例，
而是在\textbf{发现}最大化奖励的行为模式。

顺序缩放的一个微妙之处是\textbf{推理长度与质量}的关系。
并非越长的推理越好，
模型有时会陷入无意义的循环重复，
或者在简单问题上过度展开推理链。
DeepSeek-R1 的训练中观察到，
模型在 RL 早期会出现 “overthinking” 现象，
推理链越来越长但答案并没有变好。
通过调整奖励函数（对过长的推理施加轻微惩罚），
可以让模型学会在合适的深度停止思考。

\subsection{策略二：并行采样与 Best-of-N（Parallel Scaling）}

\textbf{并行缩放}对同一个问题生成 $N$ 个独立的回答，
然后用多数投票或奖励模型选最好的。
这相当于「三个臭皮匠赛过诸葛亮」，
即使每个回答独立看不够好，
通过聚合多个独立的回答可以得到更可靠的结果。

\subsubsection{Best-of-N 采样}

\textbf{Best-of-N 采样}是最简单的并行缩放实现：
生成 $N$ 个回答，
用奖励模型（Reward Model, RM）选择分数最高的那个。

数学上，Best-of-N 的效果可以用\textbf{序统计量}分析。
设每个回答的奖励分数 $r_i$ 独立同分布于某分布 $F(r)$。
$N$ 个样本的最大值 $r_{(N)}$ 的期望为：
\begin{equation}
  \mathbb{E}[r_{(N)}] = \int_{-\infty}^{\infty} r \cdot N \cdot F(r)^{N-1} \cdot f(r)\, dr
\end{equation}
当 $F$ 是正态分布 $\mathcal{N}(\mu, \sigma^2)$ 时，
$\mathbb{E}[r_{(N)}]$ 近似为：
\begin{equation}
  \mathbb{E}[r_{(N)}] \approx \mu + \sigma \cdot \Phi^{-1}\!\left(\frac{N}{N+1}\right)
\end{equation}
其中 $\Phi^{-1}$ 是标准正态分布的分位数函数。

这意味着 Best-of-N 的收益是\textbf{对数缩放}的：
从 $N=1$ 到 $N=4$，质量提升显著；
从 $N=16$ 到 $N=64$，提升已经很小；
到 $N=256$，几乎无法进一步提升。
计算成本线性增长但收益对数递减，
所以实践中 $N = 4$--$16$ 是最优的性价比区间。

\begin{warnbox}[Best-of-N 的陷阱：奖励黑客]
Best-of-N 的质量\textbf{严重依赖}奖励模型的质量。
如果奖励模型是训练目标的不完美代理
（这几乎总是如此），
过大的 $N$ 会导致\textbf{奖励黑客}（reward hacking），
模型生成的回答获得了高奖励分数，
但实际质量并不好，甚至更差。
这是因为大 $N$ 给了模型更多机会
发现奖励模型的弱点并加以利用。

\textbf{正则化 Best-of-N}（RBoN）~\cite{ichihara2025rbon}
通过在目标函数中加入正则项来缓解这个问题：
\begin{equation}
  x^* = \arg\max_{x \in \{x_1, \ldots, x_N\}} \left[ R(x) - \beta \cdot \mathrm{KL}(q \| p_{\text{ref}}) \right]
\end{equation}
正则项惩罚偏离参考分布过远的回答，
防止奖励黑客。
\end{warnbox}

\subsubsection{多数投票与自一致性}

\textbf{多数投票}（Majority Voting / Self-Consistency）~\cite{wang2023selfconsistency}
更适合有\textbf{可验证答案}的问题（如数学题）：
生成 $N$ 个推理过程和答案，
选择出现次数最多的答案。
直觉是：即使每次推理的中间步骤不同，
正确答案更有可能被多次「发现」。

数学分析很清晰：
如果单次回答正确的概率是 $p$，
$N$ 次独立采样中至少有一次正确的概率是：
\begin{equation}
  P(\text{at least one correct}) = 1 - (1-p)^N
\end{equation}
当 $p = 0.3$（单次正确率 30\%）时，
$N = 10$ 次采样就有 $1 - 0.7^{10} \approx 97\%$ 的概率
至少得到一个正确答案。
但从 $N = 10$ 到 $N = 100$，
概率只从 97\% 提升到 99.99\%，
\textbf{收益急剧递减}。

多数投票比 Best-of-N 更\textbf{鲁棒}，
因为它不依赖可能有偏的奖励模型，
它依赖的是「正确答案比错误答案更一致」这个假设。
但它只适用于答案可以\textbf{精确比较}的场景
（数学、选择题、代码输出），
不适用于开放式文本生成。

\subsubsection{奖励模型引导搜索}

更高级的并行缩放策略结合了
奖励模型的评估能力和搜索算法。

\textbf{Weighted majority voting}：
不是简单地数每个答案出现的次数，
而是用奖励模型对每个推理链打分，
然后按分数加权投票。
这结合了自一致性（答案频率）
和质量评估（奖励分数）的优点。

\textbf{LE-MCTS}（Language model Ensemble with MCTS）~\cite{park2025lemcts}：
将多个 LLM 的推理步骤视为
MCTS 中不同的行动选择。
在每一步推理中，选择哪个 LLM 来生成下一步
由 UCB（Upper Confidence Bound）策略决定，
由 PRM 提供每一步的评分。
这实现了\textbf{步骤级别}的模型集成，
比简单的输出级别集成更加精细。

并行缩放的优点是实现极其简单、可线性扩展、
而且每个独立样本可以并行生成，延迟不增加（只增加吞吐需求）。
缺点是效率低，大部分生成是浪费的，
而且对于需要深度推理的问题，
100 次浅层思考可能不如 1 次深度思考。

\subsection{策略三：搜索（Search / Tree of Thought）}

\textbf{搜索}将推理过程视为一棵\textbf{树}，
每个节点是一个中间推理步骤，
每条边是一步推理。
通过系统性地探索多条推理路径，
可以避免单条路径的死胡同。

\textbf{Tree of Thought}（ToT）
是最直观的实现：
在推理的每个关键步骤，
生成多个候选的下一步，
评估每个候选的质量，
选择最有希望的继续展开。

\textbf{MCTS}（Monte Carlo Tree Search）
将 AlphaGo 的搜索算法引入语言推理：
从根节点（问题）出发，
通过选择（Selection）、扩展（Expansion）、
模拟（Simulation）、回传（Backpropagation）
四个步骤反复迭代，
逐步构建推理树并找到最优路径。

搜索的关键组件是\textbf{评估函数}：
如何判断一个中间推理步骤的质量？
这里引入了 \textbf{Process Reward Model}（PRM）的概念：
与传统的 Outcome Reward Model（ORM，只评分最终答案）不同，
PRM 对\textbf{每个推理步骤}打分。
这提供了更细粒度的信号，
不仅知道最终答案对不对，
还知道从哪一步开始出了问题。
PRM 使得搜索可以\textbf{及早剪枝}错误的推理路径，
大幅提高搜索效率。

\begin{whybox}[ORM vs PRM：粗粒度 vs 细粒度奖励]
考虑一个数学问题的 5 步推理过程，
其中第 3 步出了错误。

\textbf{ORM}只看最终答案：答案错了 → 整个推理链得低分。
但 ORM 不知道前 2 步是对的，
导致在搜索中可能丢弃那些「前半部分优秀」的推理路径。

\textbf{PRM}对每一步评分：
Step 1: 0.95, Step 2: 0.92, Step 3: 0.15, Step 4: 0.10, Step 5: 0.08。
PRM 清晰地定位了错误在 Step 3，
搜索算法可以保留 Step 1--2 的推理，
只从 Step 3 开始尝试新的方向。

PRM 的训练更困难，
需要\textbf{步骤级别}的标注，
而不是简单的「答案对不对」。
但其提供的细粒度信号
使得搜索效率提升数个数量级。
Snell et al.\ 的实验表明，
PRM 引导的 beam search 在相同计算预算下
比 ORM 引导的 Best-of-N 高出 10--20\% 的准确率。
\end{whybox}

\subsubsection{PRM 的训练方法论}

训练一个高质量的 PRM 是推理时缩放中
最具挑战性的技术问题之一。
核心难题在于：\textbf{如何获取步骤级别的训练标签？}

\textbf{方法一：人类标注}（最高质量，最高成本）。
OpenAI 的 PRM800K~\cite{lightman2024lets}
是目前最知名的人类标注 PRM 数据集。
标注者需要对每一步推理判断：
“正确”“中性”或“错误”。
人类标注的准确率极高，
但成本惊人，
800K 步骤的标注仅覆盖了 12K 道数学题，
每道题的标注成本在 \$5--\$20 之间。
更重要的是，标注者需要\textbf{理解数学推理}，
这严重限制了标注规模的扩展，
不能像文本分类那样外包给非专业标注员。

\textbf{方法二：MCTS Rollout 自动标注}（高质量，高计算成本）。
Math-Shepherd~\cite{wang2024mathshepherd} 开创的方法：
对于推理链中的每一步 $s_i$，
从该步出发用 LLM 做 $K$ 次 rollout 到最终答案，
计算正确率 $p_i = \text{correct}(s_i) / K$。
$p_i$ 作为步骤 $s_i$ 的“过程奖励”标签。

这种方法的巧妙之处在于：
它将\textbf{步骤级别}的评估问题
转化为了\textbf{结果级别}的验证问题，
不需要人类判断“这一步推理是否正确”，
只需要验证“从这一步出发能否得到正确答案”。
后者可以完全自动化（数学题可以数值验证，代码题可以运行测试）。

计算成本是主要的限制：
对于一条 $S$ 步的推理链，
每步做 $K$ 次 rollout，每次 rollout 生成 $L$ 个 token。
总 token 生成量为 $S \times K \times L$，
如果 $S = 10$, $K = 64$, $L = 200$，
一道题的标注成本是 $128,000$ 个 token，
相当于一次\textbf{深度推理}的成本。
规模化到 100K 道题就需要约 $10^{10}$ token 的生成量，
相当于一次小规模预训练的计算量。

\textbf{方法三：ORM-to-PRM 蒸馏}（中等质量，低计算成本）。
使用已有的 ORM（只评估最终答案）
来为每一步生成\textbf{伪标签}。
具体做法是：
对于推理链的每一步 $s_i$，
移除 $s_i$ 之后的所有步骤，
让 ORM 对“截断到 $s_i$ 的推理链”打分。
如果 ORM 给出高分，
说明前 $i$ 步“看起来在正确的方向上”，
反之则可能出了问题。

这种方法的质量取决于 ORM 的能力，
如果 ORM 本身不够好（无法从部分推理链中
准确判断最终答案的可能性），
生成的伪标签会有大量噪声。
但它的计算成本极低（每步只需要一次 ORM 前向传播），
适合作为大规模 PRM 训练的\textbf{初始数据源}，
后续可以通过少量人类标注或 MCTS rollout 数据来精调。

\textbf{PRM 的粒度选择}也至关重要。
“一步推理”的定义可以很细
（每个句子一步）或很粗（每个主要推理阶段一步）。
过细的粒度（句子级别）提供了最精细的反馈信号，
但标注成本指数增长，
且 PRM 需要学习区分“微小的措辞变化是否影响正确性”：
这对 PRM 的能力要求极高。
过粗的粒度（阶段级别）标注成本低，
但信号太粗，
一个“错误阶段”中可能包含 5 个正确步骤和 1 个关键错误，
PRM 无法定位到具体的错误。
实践中，\textbf{自然推理步骤}（以换行或“Therefore”等标记分割）
是最常用的粒度：
既与 LLM 的生成结构对齐，
又提供了足够细的反馈分辨率。

\textbf{AlphaProof} 和 \textbf{AlphaGeometry 2}
是搜索策略的极致展示：
在 2024 年国际数学奥林匹克（IMO）中，
AlphaProof 与 AlphaGeometry 2 联合解出了 6 道题中的 4 道
（其中 AlphaProof 解决了代数和数论题，
AlphaGeometry 2 解决了几何题），
达到了银牌水平。
AlphaProof 使用 Gemini 模型将自然语言转换为 Lean 4 形式化语言，
再通过类似 AlphaZero 的强化学习方法进行证明搜索。
每一步都可以被\textbf{机器验证}（proof checker），
消除了幻觉的可能性。

\textbf{Beam search over reasoning chains}
是一种更轻量的搜索方法：
类似于传统的 beam search，
维护 $B$ 条最有希望的推理链，
每步扩展每条链，保留总分最高的 $B$ 条。
PRM 提供每一步的评分。
这比完整的 MCTS 计算量更小，
但仍然比纯顺序推理有显著的质量提升。

\subsubsection{搜索策略的深入对比}

不同搜索策略之间的选择
本质上是在\textbf{搜索广度}和\textbf{搜索深度}之间做权衡。
理解每种策略的计算-质量权衡曲线
是选择最优推理策略的基础。

\textbf{Best-of-N（朴素并行搜索）}
是搜索的最简单形式，完全不在“步骤”层面做选择，
只在完整输出层面做选择。
它的计算成本是 $N \times C_{\text{single}}$
（$N$ 条独立推理链的总生成成本）。
优点是实现极其简单（只需要能采样和评分），
且 $N$ 条链可以\textbf{完全并行}，延迟不增加。
但 Best-of-N 对计算的利用效率很低：
$N$ 条链中可能有多条探索了相似的错误路径，
没有任何机制来“从错误中学习”或“避免重复尝试”。
它适合低难度到中难度的问题，
在高难度问题上很快就“撞到天花板”，
因为如果单次正确率 $p$ 很低（如 5\%），
即使 $N = 100$ 也只有 $1 - 0.95^{100} \approx 99.4\%$ 的概率
至少得到一个正确答案，
但前提是你能\textbf{识别}出哪个是正确的：
这需要一个高质量的 ORM 或可验证的答案格式。

\textbf{Beam search}
在搜索的每一步都做选择，保留最好的 $B$ 条路径。
计算成本约为 $B \times S \times C_{\text{step}}$
（$B$ 条束、$S$ 步、每步的生成和评分成本）。
Beam search 的关键优势是它在推理过程中\textbf{及时剪枝}，
一旦某条路径的 PRM 分数明显低于其他路径，
立刻丢弃，将计算资源分配给更有希望的路径。
这使得 beam search 在固定计算预算下
比 Best-of-N 更有效地探索搜索空间。

但 beam search 有一个重要的弱点：\textbf{缺乏回溯}。
一旦在某一步选错了方向
（所有 $B$ 条束都选择了错误的推理方向），
后续步骤无法纠正，
因为 beam search 只\textbf{向前}扩展，不会回到之前的步骤重来。
这就是为什么在极高难度的问题上，
beam search 的表现不如 MCTS。

\textbf{MCTS（蒙特卡洛树搜索）}
是这三种策略中最强大也最昂贵的。
MCTS 的核心循环包含四步：
\begin{enumerate}[noitemsep]
  \item \textbf{选择}（Selection）：
        从根节点开始，用 UCB（Upper Confidence Bound）公式
        选择最有“潜力”的子节点（平衡“探索未知”和“利用已知”）：
        \begin{equation}
          \text{UCB}(s) = \bar{r}(s) + c \cdot \sqrt{\frac{\ln N_{\text{parent}}}{N_s}}
        \end{equation}
        其中 $\bar{r}(s)$ 是该节点的平均奖励，$N_s$ 是访问次数，
        $c$ 是探索系数。
  \item \textbf{扩展}（Expansion）：
        到达叶子节点后，用 LLM 生成一步推理，创建新的子节点。
  \item \textbf{模拟}（Simulation / Rollout）：
        从新节点出发，用快速策略（如 LLM 的 greedy decode）
        生成一条完整的推理链到达终点。
  \item \textbf{回传}（Backpropagation）：
        将终点的奖励（ORM 或验证结果）
        \textbf{反向传播}到路径上的所有祖先节点，
        更新它们的 $\bar{r}$ 和 $N$。
\end{enumerate}

MCTS 的关键能力是\textbf{回溯}：
当一个方向被证明是死胡同后，
MCTS 会回到更早的分支点尝试新方向。
UCB 公式保证了探索-利用的平衡：
不会无限地在一个方向上深挖，
也不会浅尝辄止地频繁跳转。

\begin{whybox}[MCTS 在语言推理中的挑战]
MCTS 在围棋（AlphaGo）中取得了巨大成功，
但将其应用于语言推理面临几个独特的挑战：

（1）\textbf{搜索空间的维度}。
围棋的每步有 $\sim$250 种合法落子；
语言推理的每步可能有 $\sim$50,000 种 token 选择。
这使得直接在 token 级别做 MCTS 几乎不可行，
需要在更粗的粒度（如“推理步骤”或“句子”级别）
定义搜索树的节点。

（2）\textbf{评估函数的可靠性}。
围棋有完美的评估函数（最终胜负）；
语言推理的 PRM 是一个不完美的近似：
它可能在某些推理模式上系统性地“看走眼”。
MCTS 对评估函数的可靠性极为敏感：
如果 PRM 在某类推理步骤上系统性地给高分
（即使该步骤实际上是错的），
MCTS 会被\textbf{误导}到错误的搜索方向。

（3）\textbf{rollout 的成本}。
围棋的 rollout 可以用极快的随机策略完成（毫秒级）；
语言推理的 rollout 需要 LLM 生成完整的推理链
（可能数千 token，耗时数秒）。
这使得 MCTS 的每次迭代成本比围棋高出 $10^3$--$10^4$ 倍。
AlphaProof 通过将自然语言转换为 Lean 4 形式化语言来缓解这一问题，
形式化验证可以在毫秒内判断一个推理步骤是否正确，
消除了对不完美 PRM 的依赖。
\end{whybox}

\begin{table}[H]
\centering\small
\renewcommand{\arraystretch}{1.3}
\caption{推理时搜索策略的计算-质量权衡}
\begin{tabularx}{\linewidth}{l c c c X}
\toprule
\rowcolor{figLightGray}
\textbf{策略} & \textbf{计算成本} & \textbf{回溯能力} & \textbf{需要 PRM} & \textbf{适用场景} \\
\midrule
Best-of-$N$ & $N \times 1\times$ & 无 & 推荐（ORM 可用） & 低-中难度，可精确验证的任务 \\
Beam search & $B \times S \times 1\times$ & 无 & 是（步骤级） & 中难度，结构化推理 \\
MCTS & $I \times R \times 1\times$ & 是 & 是（强依赖） & 高-极高难度，形式化验证可用 \\
\bottomrule
\end{tabularx}
\end{table}

\subsection{三种策略的对比}

\begin{keybox}[推理时计算策略对比]
\begin{itemize}[noitemsep]
  \item \textbf{顺序缩放}：一次深度思考。
        优势，推理深度高，适合需要多步推导的问题。
        劣势，单次推理可能走进死胡同。
  \item \textbf{并行缩放}：$N$ 次独立思考。
        优势，实现简单，可并行，延迟不增加。
        劣势，效率低，收益递减，可能遭受奖励黑客。
  \item \textbf{搜索}：系统性探索推理树。
        优势，可以回溯，避免死胡同，理论上最优。
        劣势，需要 PRM，计算量大，实现复杂。
\end{itemize}
经验法则：低难度问题用并行缩放，
中高难度用顺序缩放，
极高难度（数学竞赛、形式化证明）用搜索。
\end{keybox}

\subsection{混合缩放（Hybrid Scaling）}

\textbf{混合缩放}结合并行和顺序，
并行生成多个推理链，
每条链内部是顺序推理，最后取最优。
这种组合策略的直觉很清晰：
顺序缩放提供推理\textbf{深度}（每条链内的多步推导），
并行缩放提供推理\textbf{多样性}（不同链可能探索不同的解题路径）。
两者结合可以避免单一策略的弱点：
纯顺序可能走进死胡同而无法恢复，
纯并行的每条链可能都太浅而无法触及正确答案。

ThreadWeaver（2025 年提出）是混合缩放的代表性工作。
其核心创新是将推理链的生成\textbf{流水线化}，
在一条推理链的 decode 阶段进行的同时，
开始另一条推理链的 prefill，
使得多条链的计算在时间维度上重叠。
这实现了 $1.53\times$ 的延迟降低同时保持准确率，
关键在于多条推理链\textbf{不需要串行执行}，
可以在 GPU 的不同计算流上并发进行。

实践中最有效的配置通常是混合的：
让推理模型生成 $N$（通常 4--16）个独立的长推理链，
然后用多数投票或奖励模型选最优。
这同时利用了「深度思考」和「多样性」的优势。
具体的 $N$ 值选择取决于任务特性和计算预算：
对于数学竞赛题（答案可精确验证），
$N = 16$ 配合多数投票是常见选择：
单次正确率 50\% 的问题，16 次采样后多数投票的正确率
可以提升到 90\% 以上。
对于开放式推理任务（答案无法精确比较），
$N = 4$--$8$ 配合 PRM 排序是更务实的选择：
过大的 $N$ 会增加奖励黑客的风险。

混合缩放的一个微妙问题是
\textbf{多样性与质量的权衡}。
如果所有推理链使用相同的采样温度，
生成的推理路径可能高度相似，
降低了并行采样的边际收益。
提升多样性的常用方法包括：
（1）使用较高的采样温度（$T = 0.7$--$1.0$）
以增加随机性，但过高的温度会降低单条链的质量；
（2）在不同的链中使用不同的系统 prompt 或引导策略，
鼓励模型从不同角度切入问题；
（3）使用不同的推理模型，
如一条链用 R1，另一条用 Qwen3，
利用不同模型的推理偏好获得更大的多样性。
LE-MCTS~\cite{park2025lemcts} 将这种多模型策略
推向了步骤级别的粒度，
在每一步推理中动态选择最合适的模型。

\subsection{收益递减与最优停止}

所有推理时缩放策略都面临\textbf{收益递减}（diminishing returns）。
理解收益递减的规律
是设计计算高效的推理系统的基础。

\textbf{顺序缩放的收益递减}
遵循近似对数曲线。
DeepSeek-R1 的实验数据表明：
在 AIME 数学题上，
thinking tokens 从 256 增加到 2K（$8\times$）时，
准确率从约 55\% 提升到约 75\%（+20\%）；
从 2K 增加到 16K（$8\times$）时，
准确率从 75\% 提升到约 80\%（+5\%）；
从 16K 增加到 128K（$8\times$）时，
准确率仅从 80\% 提升到约 82\%（+2\%）。
每次 $8\times$ 的计算增加，
准确率提升幅度缩小约 $4\times$，
这是典型的\textbf{对数递减}模式。

数学上，如果将准确率建模为推理计算 $C$ 的函数：
\begin{equation}
  A(C) \approx A_\infty - \frac{k}{C^\alpha}, \quad \alpha \in (0.3, 0.5)
\end{equation}
其中 $A_\infty$ 是准确率的\textbf{上限}
（由模型的知识和推理能力决定），
$k$ 是问题难度相关的常数，
$\alpha$ 是缩放指数。
这个公式暗示：
（1）准确率有一个\textbf{不可突破的天花板} $A_\infty$；
（2）接近天花板时，额外计算的边际收益趋近于零；
（3）对于不同难度的问题，天花板 $A_\infty$ 不同。

\textbf{并行缩放的收益递减}更为剧烈。
Best-of-$N$ 的准确率改善大致遵循 $\sqrt{\log N}$，
意味着从 $N = 1$ 到 $N = 4$ 的收益
几乎等于从 $N = 4$ 到 $N = 64$ 的收益。
$N > 100$ 之后的边际收益在实践中几乎不可观测。

\textbf{最优停止策略}是应对收益递减的实用工具。
核心思想是：不要总是用完预分配的计算预算：
如果模型在推理链中\textbf{早期}就产生了高置信度的答案，
应该提前终止推理，将节省的计算用于其他请求。

实现最优停止的技术方案包括：
\begin{itemize}[noitemsep]
  \item \textbf{置信度阈值}：
        在推理链的每个主要步骤后，
        用 PRM 或 ORM 评估当前答案的置信度。
        如果置信度超过阈值（如 0.95），提前返回答案。
  \item \textbf{一致性检测}：
        在并行采样中，如果前 $k$ 个样本
        已经一致地给出同一个答案（如 4/4 一致），
        提前停止采样，
        因为更多样本几乎不会改变多数投票的结果。
  \item \textbf{进展检测}：
        监控推理链中“有效进展”的速率。
        如果模型在最近 500 个 token 中
        没有引入新的推理步骤或信息
        （只是在重复或犹豫），
        可能已经陷入了无效循环，应该触发终止。
\end{itemize}

最优停止策略在生产环境中的价值不仅在于节省计算，
更重要的是它\textbf{降低了 overthinking 的风险}。
过度思考不仅浪费资源，
有时甚至会\textbf{降低}答案质量，
模型在过长的推理链中可能“自我说服”采纳一个错误方向，
导致最终答案反而不如短推理链的直觉回答。
S1~\cite{muennighoff2025s1} 的研究发现，
对于某些中等难度的数学题，
最优的 thinking token 数量是 500--1000，
超过 2000 后准确率反而\textbf{下降}了 2--5\%。
budget forcing 和 early exit 可以防止这种过度思考。


\section{DeepSeek-R1：从零开始学会推理}

DeepSeek-R1~\cite{deepseekr1} 是推理时缩放领域的里程碑。
它首次证明了\textbf{纯粹的强化学习}（无需人类推理示范）
就能让模型自发地发展出复杂的推理能力。

2025 年 9 月，R1 论文以封面文章发表于
\textit{Nature}（Vol.~645, Issue~8081, 2025 年 9 月 17 日），
成为\textbf{首篇通过 Nature 独立同行评审的 LLM 论文}。
梁文锋（Liang Wenfeng）为通讯作者。
Nature 编辑部以“AI can learn to show its workings through trial and error”为题
发表评论，确认了 R1 的核心发现：
纯 RL（无需人类标注的推理轨迹）能够激励高级推理行为的涌现，
包括自我反思（self-reflection）、
验证（verification）和动态策略调整（dynamic strategy adaptation）。
这一里程碑标志着 LLM 推理研究
从工程实践迈入了经由严格同行评审认可的科学领域。

\subsection{R1-Zero：从 RL 中涌现的推理}

在 R1 的完整训练之前，
DeepSeek 团队先进行了一个大胆的实验：
在 DeepSeek-V3-Base 上\textbf{直接}应用 GRPO RL 训练，
\textbf{不使用任何 SFT 数据}。
这个模型被称为 \textbf{DeepSeek-R1-Zero}。

R1-Zero 的训练设计极其简洁：
\begin{itemize}[noitemsep]
  \item \textbf{奖励信号}：只有两类，
        数学题答案是否正确（accuracy reward），
        和输出格式是否符合要求（format reward）。
        没有过程奖励（PRM），没有人类偏好。
  \item \textbf{训练数据}：数学、代码和逻辑推理题，
        全部有\textbf{可验证}的正确答案。
  \item \textbf{优化算法}：GRPO，
        对每个问题生成一组回答，
        根据组内相对排名计算优势函数，
        不需要单独的价值网络或奖励模型。
\end{itemize}

\subsection{“Aha Moment”：模型的顿悟时刻}

R1-Zero 训练过程中出现了令研究者震惊的现象。
在 RL 训练到一定阶段后，
模型自发地开始展现出一系列
\textbf{从未在训练数据中出现过}的推理行为：

\begin{corebox}[DeepSeek-R1 的 “Aha Moment”]
在训练中期的某次评估中，
模型在解数学题时写下了这样的推理过程：

\textit{“Wait, wait. Wait. That's an aha moment.
I can actually solve this by re-reading the problem
and approaching it differently...
Let me reconsider what I assumed earlier...”}

这不是人类编写的推理模板：
模型自己「发明」了自我反思和重新审视的策略。
这个时刻被 DeepSeek 团队称为 “aha moment”，
因为它标志着模型从简单的模式匹配
跃迁到了\textbf{元认知}层面的推理。
\end{corebox}

具体而言，R1-Zero 自发发展出了以下推理策略：

\textbf{回溯}（Backtracking）：
模型在推理链中发现前面的方向不对时，
主动放弃当前路径，回退到更早的步骤重新推导。
这是一种典型的搜索行为：
模型学会了在推理树中「剪枝」和「回溯」。

\textbf{自我验证}（Self-verification）：
模型在得出中间结论后，
会主动检验该结论的正确性，
例如将答案代回原方程验证，
或用不同方法重新推导检查一致性。
这模拟了数学家的证明检查过程。

\textbf{分类讨论}（Case analysis）：
面对多种可能情况的问题，
模型学会了系统性地枚举每种情况，
逐一分析，最后综合得出结论。

\textbf{反证法}（Proof by contradiction）：
在某些逻辑推理问题中，
模型会假设结论不成立，
推导出矛盾，从而证明结论成立。

这些策略的涌现具有深远的意义：
它表明复杂的推理能力\textbf{不需要}显式的人类教学，
可以通过适当的 RL 训练从模型中自然涌现。
这与 AlphaGo Zero 的启示一致，
在有可验证奖励信号的领域，
RL 能够发现人类未曾想到的策略。

\textbf{涌现的可复现性}
是 R1-Zero 实验中一个被低估的重要发现。
DeepSeek 团队报告称，
在不同的随机种子下重复 R1-Zero 的训练，
推理行为的涌现\textbf{时间点}有所不同
（从 500 步到 2000 步不等），
但涌现的\textbf{行为类型}是一致的：
回溯、自我验证、分类讨论
在所有实验中都出现了。
这表明这些推理策略不是特定随机种子的偶然产物，
而是 RL 优化在可验证奖励下的\textbf{必然涌现}。

这个发现有重要的理论含义。
它暗示推理策略空间中存在一组
\textbf{“吸引子”}（attractors），
无论从哪个初始点出发，
RL 优化都会收敛到这些策略。
回溯、验证、分类讨论之所以是“吸引子”，
可能是因为它们在最大化奖励方面
具有\textbf{结构性优势}：
回溯允许模型从错误中恢复
（增加了正确答案的概率），
验证减少了虚假答案
（减少了不必要的奖励损失），
分类讨论确保了对各种情况的覆盖
（在有多种可能情况的问题中尤为有效）。

\textbf{推理策略的“文化”传播}
是 R1 发布后的一个有趣的社区现象。
开源的 R1 推理链成为了大量后续模型的\textbf{训练数据}，
社区中的蒸馏模型、微调模型
都在不同程度上“继承”了 R1 的推理风格。
这导致了推理模型生态中的
一种“文化同质化”，
几乎所有中文推理模型的思考过程
都带有 R1 特征性的“Wait, let me reconsider...”风格。
这既是 R1 影响力的证明，
也引发了一个\textbf{多样性关切}：
如果所有推理模型都“像 R1 一样思考”，
推理策略的多样性可能下降，
限制了通过集成不同模型来提升推理质量的空间。

\subsection{R1 的完整训练流程}

R1-Zero 虽然展现了令人兴奋的推理涌现，
但存在两个严重的问题：
\begin{itemize}[noitemsep]
  \item \textbf{可读性差}：输出混合多种语言，
        推理链混乱难读，对用户不友好。
  \item \textbf{非推理任务退化}：
        在日常对话、翻译等不需要推理的任务上
        表现明显下降。
\end{itemize}

DeepSeek-R1 的完整训练分为四个阶段：

\textbf{阶段一：冷启动 SFT}。
收集少量高质量的推理示范数据（数千条），
对基础模型做 SFT。
这一步不是为了教模型推理：
而是为了稳定后续 RL 训练的起点，
避免 R1-Zero 中的格式混乱问题。

\textbf{阶段二：推理导向 RL}。
在冷启动模型上做 GRPO RL 训练，
使用数学、代码和逻辑题的可验证奖励。
这是推理能力形成的核心阶段。

\textbf{阶段三：拒绝采样 + 通用 SFT}。
用阶段二的模型生成大量推理链，
通过拒绝采样保留高质量的输出（约 800K 条），
混合通用的 SFT 数据（日常对话、翻译等），
做一轮全面的 SFT。
这一步恢复了模型在非推理任务上的能力。

\textbf{阶段四：全场景 RL}。
在所有类型的任务上做最终的 RL 训练，
使用混合的奖励信号
（可验证任务用规则奖励，开放式任务用奖励模型），
进一步提升整体质量和对齐度。

这四个阶段的设计体现了一个重要的工程原则：
\textbf{能力的分阶段注入}。
如果直接在 base 模型上做全场景 RL
（跳过冷启动和推理导向 RL），
模型难以同时学会“推理”和“遵循格式”，
两个目标互相干扰，导致训练不稳定。
将训练分为“先学格式（冷启动）→ 再学推理（RL）
→ 恢复通用能力（SFT）→ 最终对齐（全场景 RL）”，
每个阶段都有明确的目标和稳定的梯度信号。

一个经常被忽视的细节是
\textbf{阶段三的数据配比}对最终模型的影响。
如果推理蒸馏数据占比过高（如 90\%），
模型在推理任务上表现极佳
但在日常对话中显得“过度推理”，
回答“你好”时也会展开一段不必要的推理链。
如果通用 SFT 数据占比过高（如 90\%），
模型的推理能力被“稀释”，
R1 团队发现 60:40（推理:通用）的配比
在推理能力和通用能力之间取得了最佳平衡。

R1 训练流程的另一个关键洞察是
\textbf{数据的“自举”效应}。
阶段二的 RL 模型虽然格式不够好，
但推理质量已经相当高。
用这个模型做\textbf{拒绝采样}生成的推理链
成为了阶段三 SFT 的训练数据，
模型用自己的输出来训练自己的下一个版本。
这种自举式的数据生成
使得 R1 不依赖外部的人类推理标注，
它“自己创造了自己的训练数据”。

\subsection{R1 vs o1/o3：两条不同的路径}

DeepSeek-R1 和 OpenAI 的 o1/o3 系列
代表了推理时缩放的两条不同路径：

\begin{table}[H]
\centering\small
\renewcommand{\arraystretch}{1.3}
\caption{DeepSeek-R1 vs OpenAI o 系列对比}
\begin{tabularx}{\linewidth}{l X X}
\toprule
\rowcolor{figLightGray}
\textbf{维度} & \textbf{DeepSeek-R1} & \textbf{OpenAI o1/o3} \\
\midrule
\textbf{核心方法} & GRPO RL + 可验证奖励。不需要单独的奖励模型。 & RL + 详细的 PRM。使用复杂的奖励建模系统。 \\
\textbf{推理过程} & 完全可见。用户可以看到完整的 thinking tokens。 & 隐藏。用户只看到最终答案（部分摘要可见）。 \\
\textbf{开源} & 模型权重完全开源，训练方法论在论文中详细公开。 & 完全闭源，仅通过 API 提供服务。 \\
\textbf{基础架构} & 671B MoE（37B 活跃参数），基于 DeepSeek-V3。 & 未公开。推测基于 GPT-4 级别架构。 \\
\textbf{蒸馏} & 提供 1.5B 到 70B 的完整蒸馏系列。 & 无官方蒸馏，o3-mini 作为轻量版本。 \\
\textbf{性能（AIME 2024）} & 79.8\% & o1: 83.3\%, o3-mini (high): 87.3\% \\
\bottomrule
\end{tabularx}
\end{table}

o3-mini 是一个值得关注的案例：
它在多个数学基准（如 AIME）上\textbf{超过}了 o1，
同时成本只有 o1 的 $1/14$（\$1.10/\$4.40 vs \$15/\$60 per M tokens）。
这证明了推理时缩放中
\textbf{算法效率的提升}可以同时提高质量和降低成本。

2025 年 4 月，OpenAI 发布了 o3 和 o4-mini，
标志着推理模型的又一次飞跃。
o3 和 o4-mini 引入了\textbf{工具使用}能力，
它们可以在推理过程中自主决定
搜索网页、执行代码、分析图像，
将工具调用和推理无缝结合。
这代表了从「纯思考」到「思考 + 行动」的范式转变。

\subsection{“Learning to Think”：训练模型学会思考}

o1 和 R1 代表的“学会思考”范式
与传统的 Chain-of-Thought prompting 有本质区别。
传统 CoT 是\textbf{外部注入}的：
人类通过 prompt 工程告诉模型“请一步步思考”，
模型被动地遵循指令。
o1/R1 的推理能力是\textbf{内在习得}的：
通过大规模 RL 训练，
模型\textbf{自己发现}了“思考”比“直接回答”能获得更高奖励，
因此主动选择生成长推理链。

训练模型“学会思考”的核心技术挑战在于\textbf{信用分配}（credit assignment）。
一条 10,000 token 的推理链中，
可能只有少数几百个 token 是“关键推理步骤”，
其余是冗余的重复、犹豫或无效的探索。
RL 的奖励信号只在最终答案处给出，
模型如何知道“哪些 thinking token 是有效的”？

GRPO 的巧妙之处在于它通过\textbf{组内比较}
来间接解决信用分配问题。
对同一个问题生成一组（如 8 条）不同的推理链，
正确的推理链获得正向奖励，
错误的获得负向奖励。
通过比较正确链和错误链之间的\textbf{差异}，
梯度信号隐式地指向了
“导致正确和错误差异的关键 token”。
这不是完美的信用分配：
正确链中的冗余 token 也会获得不应得的正向梯度，
但在足够大的训练规模下，
统计上有效的推理模式会被\textbf{一致地强化}，
而随机的冗余内容会被\textbf{平均化}掉。

\textbf{推理长度的涌现性增长}
是训练过程中最令人惊讶的现象之一。
DeepSeek 报告称，R1-Zero 的平均 thinking token 数量
在 RL 训练过程中\textbf{自发增长}：
从训练初期的约 200 tokens
增长到训练后期的约 3,000 tokens。
模型“学会了”投入更多计算来获取更高奖励，
这是一种\textbf{学习到的计算-质量权衡}。
更深层次地看，推理链长度的增长
反映了模型内部推理策略复杂度的增长，
从简单的“直接计算”进化为包含
回溯、验证、分类讨论的多层策略。

\textbf{训练与推理计算的互补关系}
可以用一个\textbf{等性能曲面}来刻画。
设训练计算为 $C_{\text{train}}$，
推理计算为 $C_{\text{infer}}$，
目标准确率为 $A$。
等性能曲面 $C_{\text{train}} \times C_{\text{infer}} = f(A)$
描述了达到准确率 $A$ 所需的\textbf{最小总计算量}。
曲面上的每一点代表一种不同的“训练 vs 推理”计算分配。

关键发现是这个曲面\textbf{不是线性的}：
\begin{itemize}[noitemsep]
  \item 对于简单问题（$A < 0.7$），
        等性能曲面近似 $C_{\text{train}}^{0.5} + C_{\text{infer}}^{0.1} = \text{const}$，
        训练计算的边际效益远高于推理计算，
        应该优先投入训练。
  \item 对于中难度问题（$0.7 < A < 0.95$），
        等性能曲面变得更“对称”，
        训练和推理计算的边际效益接近，
        最优策略是两者\textbf{均衡分配}。
  \item 对于极难问题（$A > 0.95$），
        推理计算的边际效益反超训练计算，
        因为极难问题需要的不是更多知识，
        而是更深的搜索和验证。
\end{itemize}

2026 年的前沿研究正在尝试\textbf{自动化}这种分配决策，
根据目标任务分布和硬件预算，
自动确定训练模型的最优大小
和推理时的最优计算分配。
这是一个\textbf{元优化}问题，
答案取决于任务的难度分布、硬件的算力-成本比
和延迟 SLO 约束。

\subsection{推理基准的军备竞赛（2025--2026）}

DeepSeek-R1 发表后，推理模型的竞争进入白热化阶段。
2025 年末至 2026 年初，多个模型在传统推理基准上
取得了接近天花板的成绩：

\begin{table}[H]
\centering\small
\renewcommand{\arraystretch}{1.3}
\caption{推理模型基准成绩（2025 年末至 2026 年初）}
\begin{tabularx}{\linewidth}{l X r}
\toprule
\rowcolor{figLightGray}
\textbf{模型} & \textbf{关键成绩} & \textbf{发布时间} \\
\midrule
\textbf{DeepSeek V3.2-Speciale} &
  IMO 金牌（35/42）、ICPC 第 2 名（492/600）、IOI 第 10 名、AIME 96.0\%（超越 GPT-5-High 94.6\%） &
  2025-12 \\
\textbf{Doubao Seed 2.0 Pro} &
  AIME 2025 98.3\%、Codeforces Rating 3020 &
  2026-02 \\
\textbf{Step 3.5 Flash} &
  AIME 97.3\%，仅 11B 活跃参数（196B 总参数 MoE） &
  2026-02 \\
\textbf{Kimi K2.5} &
  AIME 96.1\%、HMMT 95.4\%、GPQA Diamond 87.6\% &
  2026-01 \\
\textbf{Qwen3-Max-Thinking} &
  AIME/HMMT 部分基准宣称 100\% &
  2025-09+ \\
\bottomrule
\end{tabularx}
\end{table}

值得注意的是 \textbf{DeepSeek V3.2-Speciale} 的竞赛表现：
IMO 金牌（35/42 分）意味着模型的数学推理能力
已经达到了国际数学奥林匹克顶尖选手的水平，
而 ICPC 第 2 名和 IOI 第 10 名
则证明了这种推理能力可以迁移到编程竞赛领域。
\textbf{Step 3.5 Flash} 尤其引人关注，
它仅用 11B 活跃参数就达到了 AIME 97.3\%，
证明了小型 MoE 架构配合高质量推理训练
可以在效率和性能之间取得惊人的平衡。

\begin{warnbox}[基准饱和与评估范式转移]
多个模型在 AIME 上突破 96\% 甚至宣称 100\%，
这种\textbf{分数向天花板收敛}的现象
表明传统推理基准正在快速饱和。
当区分度消失时，基准就失去了评估价值。

下一代有意义的评估正在转向更接近真实世界的任务：
\textbf{SWE-bench Pro}（1,865 个多语言真实代码修复任务，
当前最高约 44\%）衡量的是工程推理能力；
\textbf{BrowseComp}（Web agent 任务）
衡量的是模型在开放环境中的规划和执行能力。
这些基准的天花板远高于 AIME，
更能区分不同模型的实际推理水平。
\end{warnbox}


\section{验证方法：如何知道推理是否正确}

推理时缩放的效果高度依赖于
对推理结果的\textbf{验证能力}。
没有可靠的验证，更多的计算只会产生更多的“噪声”。
不同的验证方法适用于不同类型的任务。

\subsection{自一致性验证（Self-Consistency）}

\textbf{自一致性}~\cite{wang2023selfconsistency}
是最简单也最广泛适用的验证方法。
核心假设是：\textbf{正确答案比错误答案更一致}，
同一个正确的数学答案可以通过多条不同的推理路径得到，
而错误答案往往各不相同。

自一致性的工作流程：
生成 $N$ 条独立的推理链，
提取每条链的最终答案，
选择出现频率最高的答案。
它不需要任何额外的模型（ORM 或 PRM），
也不需要答案是数值型的：
只要答案可以\textbf{精确比较}即可
（数学答案、选择题选项、代码的执行输出）。

自一致性的理论基础可以用\textbf{Condorcet 陪审团定理}类比。
如果每条推理链独立地以概率 $p > 0.5$ 得到正确答案，
随着 $N$ 增大，多数投票的正确率趋近于 1。
但关键假设是“独立”：
如果模型对某类错误有\textbf{系统性偏差}
（例如总是犯同一种计算错误），
所有链可能一致地给出同一个错误答案，
多数投票反而\textbf{加强}了错误。

\begin{warnbox}[自一致性的隐含假设]
自一致性在以下条件下可能\textbf{失效}：
\begin{itemize}[noitemsep]
  \item \textbf{系统性偏差}：
        模型对某个问题有错误的“直觉”（如 “9.11 > 9.9”），
        所有推理链可能都从错误的前提出发。
  \item \textbf{低基准正确率}：
        如果 $p < 0.5$，多数投票会选择错误答案。
        当 $p = 0.2$（单次正确率只有 20\%），
        $N = 5$ 时多数投票的正确率仅约 6\%，
        \textbf{比不投票还差}。
  \item \textbf{非唯一正确答案}：
        如果一道题有多个等价但文本不同的正确答案
        （如 “$x = 2$” 和 “$x = 4/2$”），
        正确答案的“选票”被分散到多个变体上，
        单一的错误答案反而可能获得最多票数。
\end{itemize}
\end{warnbox}

\subsection{形式化验证（Formal Verification）}

对于数学和代码领域，
\textbf{形式化验证}提供了最强的正确性保证，
验证结果不是“可能正确”，而是“证明正确”。

\textbf{代码验证}最为成熟：
模型生成的代码可以直接在沙箱中运行，
与预期输出比较。
这提供了\textbf{完美的验证}，
不存在误报或漏报
（假设测试用例覆盖了所有边界条件）。
在 HumanEval 和 LiveCodeBench 等代码基准上，
代码执行验证使得 Best-of-N + 执行验证
成为一种极其有效的推理时缩放策略：
$N = 64$ 次生成中只要有一个通过所有测试，就算成功。

\textbf{数学形式化验证}更为强大但也更困难。
AlphaProof 使用 Lean 4 证明助手，
模型生成的每一步数学推理
都被翻译为形式化证明步骤，
Lean 的类型检查器自动验证每步的逻辑正确性。
这消除了“幻觉”的可能性，
如果 Lean 接受了一个证明，它在数学上\textbf{一定}是正确的。
但代价是模型必须能生成 Lean 语法的证明
（这比生成自然语言的数学推理困难得多），
且 Lean 的表达能力虽然在理论上是完备的，
但许多直觉性的数学推理
在形式化过程中会变得极其冗长。

\textbf{计算器验证}是一种轻量级的“部分形式化”。
模型在推理过程中生成计算表达式，
由外部计算器精确求值。
这不能验证推理的\textbf{逻辑}是否正确
（“从前提 A 能否推出结论 B”），
但能保证\textbf{算术}是正确的
（“$2^{32} = 4294967296$ 而不是其他数字”）。
在实践中，算术错误是 LLM 推理中最常见的失败模式之一，
特别是在多步计算（如概率计算、组合计数）中，
每步 1--5\% 的算术错误率会在 10 步之后
累积到 10--40\% 的总错误率。
计算器验证以极低的成本消除了这类错误。

\subsection{共识验证与交叉验证}

\textbf{多模型共识}（Multi-model consensus）
利用不同模型之间的“独立性”来增强验证可靠性。
核心思想是：
不同的 LLM（R1、o3、Gemini、Claude）
在训练数据和架构上各不相同，
它们的\textbf{错误模式}也不同。
当多个不同的模型独立地得到相同的答案时，
该答案正确的概率远高于单一模型的置信度。

LE-MCTS~\cite{park2025lemcts} 将多模型共识
推向了步骤级别的粒度，
在推理的每一步，使用不同的 LLM 生成候选，
通过 UCB 策略选择最优的下一步。
这不仅利用了不同模型的“知识互补”，
还利用了它们的“错误不相关性”，
即使每个模型在某步推理上有 20\% 的错误率，
三个独立模型都犯同一个错误的概率
只有 $0.2^3 = 0.8\%$。

\textbf{自我交叉验证}（Self-cross-verification）
是不需要多个模型的替代方案。
模型在得到一个答案后，
用\textbf{不同的推理路径}重新验证：
例如，用代数方法解题后，
用几何方法或数值模拟来检验结果的一致性。
R1 和 o1 的推理链中经常出现
“Let me verify this using a different approach...”
这样的自我验证行为，
这是 RL 训练自发涌现的策略，
因为“多角度验证”确实能提高最终答案的正确率，
从而获得更高的奖励。


\section{工具增强推理}

2025--2026 年，推理模型的一个重要演进方向是
将\textbf{外部工具}集成到推理过程中。
模型不再仅仅依赖内部推理，
而是可以在思考过程中调用工具来辅助推理。

\subsection{代码执行作为推理工具}

最自然的工具增强推理是\textbf{代码执行}。
在解决数学或数据分析问题时，
模型可以编写代码来精确计算中间结果，
而不是依赖自然语言推理（容易出现算术错误）。

典型的流程是：
\begin{enumerate}[noitemsep]
  \item 模型分析问题，决定使用代码来辅助计算。
  \item 在 thinking 过程中生成 Python 代码。
  \item 代码在沙箱环境中执行，返回结果。
  \item 模型将执行结果纳入后续推理链。
\end{enumerate}

这解决了 LLM 推理中最臭名昭著的弱点之一：
\textbf{算术错误}。
LLM 通过自然语言做多步算术时，
每一步都有累积误差的风险。
而代码执行保证了计算的\textbf{精确性}。
OpenAI 的 o3/o4-mini 原生支持
在推理过程中执行 Python 代码，
这使得它们在需要精确计算的数学和科学问题上
表现显著优于仅用自然语言推理的模型。

\subsection{搜索与检索增强推理}

推理过程中的\textbf{检索增强}（Retrieval-Augmented Reasoning）
让模型在思考时可以搜索外部知识。
当模型在推理中遇到不确定的事实性问题时，
它可以暂停推理、搜索相关信息、
然后基于检索到的事实继续推理。

这与传统的 RAG（Retrieval-Augmented Generation）不同：
\begin{itemize}[noitemsep]
  \item \textbf{传统 RAG}：在生成\textbf{之前}检索一次，
        将检索结果作为上下文输入模型。
  \item \textbf{推理中检索}：在推理\textbf{过程中}多次检索，
        每次检索都基于当前推理状态的需要。
        检索和推理交替进行，形成「思考 → 搜索 → 思考」的循环。
\end{itemize}

这种交互式检索在需要综合多源信息的推理任务中
效果尤为显著，
例如法律案例分析（需要检索多个先例）
或科学文献综述（需要交叉引用多篇论文）。

\subsection{形式化工具作为推理放大器}

\textbf{形式化工具}（formal tools）
代表了工具增强推理中最强大的一类。
与代码执行（验证计算结果）
和搜索引擎（提供事实信息）不同，
形式化工具可以验证\textbf{推理逻辑本身的正确性}。

\textbf{Lean 4 证明助手}在 AlphaProof 中的使用
展示了形式化工具的极限潜力。
模型生成的每一步数学推导
被自动翻译为 Lean 4 的类型论语言，
Lean 的内核（proof kernel）对每步做形式化验证。
如果某步推导存在逻辑漏洞，
Lean 会立刻报错并指出具体的类型不匹配：
模型可以根据错误信息修正推理，然后重试。
这种“生成 → 验证 → 修正”的循环
使得模型的推理过程获得了\textbf{数学级别}的正确性保障。

\textbf{SMT Solver}（Satisfiability Modulo Theories）
是另一类强大的形式化工具。
SMT 可以自动判定一阶逻辑公式的可满足性，
模型可以将推理中的关键命题
编码为 SMT 公式，让求解器自动验证。
例如，在规划任务中，
模型可以将“是否存在一个行动序列满足所有约束条件”
编码为 SMT 问题，获得确定性的“是/否”答案。
Z3~\cite{demoura2008z3}
是最广泛使用的 SMT 求解器，
其 Python 绑定使得与 LLM 的集成非常方便。

形式化工具的核心限制在于\textbf{表达性}。
并非所有推理都能翻译为形式化语言，
涉及常识推理、隐含假设、模糊概念的推理
通常无法精确形式化。
这意味着形式化工具最适合
\textbf{数学、逻辑和程序验证}等
“精确推理”领域，
对于日常推理和科学假设生成的帮助有限。

\subsection{多步 Agent 推理}

\textbf{Agent 推理}是工具增强推理的最完整形式。
推理模型不仅可以思考和使用单个工具，
还可以\textbf{规划}和\textbf{编排}多步骤的工作流。

Claude Opus 4.6 的 extended thinking
允许模型在一个统一的推理过程中
思考、调用工具、分析工具返回的结果、
再继续思考，
整个过程是一个\textbf{连贯的}推理链，
而不是分离的「思考步骤」和「工具调用步骤」。

OctoTools~\cite{lu2025octotools} 提出了一个
可扩展的 agent 推理框架。
通过\textbf{工具卡片}（tool card）机制，
模型可以在无需额外训练的情况下
集成新的工具，
每张工具卡片描述了工具的功能、输入/输出格式
和使用时机。
模型通过阅读工具卡片来决定何时使用什么工具。
实验表明，即使是最强的 LLM + 完善的 agent 工作流，
在复杂的跨论文科学推理任务上
准确率也仅有 38.78\%，
困难子集上更是低至 18.47\%：
这表明工具增强推理仍有巨大的改进空间。

\begin{keybox}[工具增强推理的层次]
\begin{itemize}[noitemsep]
  \item \textbf{Level 0，纯推理}：
        模型仅用自然语言推理，不使用任何外部工具。
        CoT、R1-Zero 属于这一层次。
  \item \textbf{Level 1，单工具推理}：
        模型在推理中使用单一工具（如代码执行器）。
        o3 的 Python 执行属于这一层次。
  \item \textbf{Level 2，多工具推理}：
        模型在推理中编排多个工具（代码、搜索、API 调用）。
        o3/o4-mini 的完整工具使用属于这一层次。
  \item \textbf{Level 3，Agent 推理}：
        模型规划多步骤工作流，
        包括子任务分解、工具编排、结果验证、
        和基于中间结果的动态调整。
        Claude 的 extended thinking + tool use 属于这一层次。
\end{itemize}
\end{keybox}


\section{推理模型的部署经济学}

推理时计算缩放正在深刻重塑 AI 的经济学。
推理需求预计将大幅超过训练需求（据行业估计可达百倍量级）。
GPT-4 等效性能的成本从 2022 年的 \$20/M tokens
降至 2026 年的 \$0.40/M tokens，
$1000\times$ 的下降，主要来自推理优化。

但推理模型也带来了新的成本挑战：
DeepSeek-R1 在回答一个数学问题时
可能生成 10,000+ 个 thinking token，
然后才给出 100 个 token 的答案。
thinking token 虽然不展示给用户，
但同样需要计算和内存资源。

\subsection{何时使用推理时缩放}

推理时缩放不是万灵药：
对于不同类型的任务，性价比差异巨大：

\textbf{高价值场景}（值得投入 10--100$\times$ 推理计算）：
\begin{itemize}[noitemsep]
  \item 数学证明和竞赛题，答案可验证，推理深度直接决定正确率。
  \item 代码 bug 修复和系统设计，错误成本高，
        多花 \$0.10 的推理费用可以避免 \$10,000 的生产事故。
  \item 法律和医疗分析，高风险决策，准确性比延迟更重要。
  \item 科学研究中的假设生成，需要深度推理和创造性思考。
\end{itemize}

\textbf{低价值场景}（标准推理即可）：
\begin{itemize}[noitemsep]
  \item 日常聊天和问答，首次回答通常就足够好。
  \item 翻译和摘要，任务结构明确，不需要深度推理。
  \item 简单的分类和提取，模型直觉足够准确。
\end{itemize}

\subsection{路由与预算管理}

这催生了\textbf{推理预算管理}的概念：
对于简单问题，限制 thinking token 的数量以节省成本；
对于困难问题，分配更多的推理预算以保证质量。

\textbf{难度路由}（Difficulty Routing）
是一种经济高效的方案：
使用一个轻量级的分类器（甚至是 prompt-based 的判断）
评估问题的难度，
将简单问题路由到标准模型（快且便宜），
将困难问题路由到推理模型（慢但准确）。
好的路由器可以将 80\% 的请求用便宜的方式处理，
只对 20\% 真正需要深度推理的问题投入额外计算。

\begin{corebox}[可配置的推理预算：行业实践]
Claude Opus 4.6 支持可配置的 thinking budget，
允许开发者通过 API 设置任意的 \texttt{budget\_tokens} 参数
来控制推理深度。
这将推理时缩放从「全有或全无」
变成了\textbf{连续可调}的维度，
开发者可以根据任务需求和成本预算
在延迟、质量和成本之间做精细的权衡。

Qwen3~\cite{qwen2025qwen3} 走了另一条路：
在单个模型中统一 thinking 和 non-thinking 模式。
用户可以根据需要在两种模式间切换——
简单问题直接回答（non-thinking），
复杂问题开启深度推理（thinking）——
不需要部署两个不同的模型。

Google Gemini 2.5/3 引入了\textbf{动态思考模式}——
模型根据问题复杂度自动调整推理深度，
无需用户手动控制。
Flash 版本提供快速推理，Pro 版本提供深度推理，
两者共享同一架构但分配不同的计算预算。
\end{corebox}

\subsection{共享推理的可能性}

对于多个用户问相似的难题（如热门的面试编程题），
可以缓存和复用推理过程。
一个用户的深度推理结果可以被后续用户共享，
避免重复计算。
这类似于 prefix caching 的思想，
但扩展到了推理链（reasoning chain）层面。

\textbf{推理链缓存的技术实现}面临几个独特挑战。
首先是\textbf{语义匹配}——
两个问题可能用不同的措辞描述了同一个数学问题，
或者一个问题是另一个的变体（改变了数字但保留了结构）。
简单的文本匹配无法捕捉这种语义相似性。
可能的方案包括：
对问题进行\textbf{规范化表示}
（例如，将“求 $\int_0^1 x^3 dx$”
和“计算 $x^3$ 从 0 到 1 的定积分”
映射到相同的表示），
或者使用嵌入模型计算问题间的语义距离，
当距离低于阈值时命中缓存。

其次是\textbf{缓存有效性}的判断。
推理链的正确性高度依赖于具体问题的细节——
一个数字的变化可能使整个推理链失效。
因此共享推理通常只适用于
\textbf{推理方法可迁移、但具体数值需要重新计算}的场景。
一种折中方案是缓存“推理模板”
（问题类型 $\to$ 推理策略 $\to$ 解题框架），
而非完整的推理链。

\textbf{隐私问题}是另一个重要考量。
如果推理链中包含了前一用户的上下文信息
（如“用户 A 在准备 Google 面试，
这道题的背景是……”），
共享给用户 B 可能构成信息泄露。
解决方案是在缓存时
对推理链进行\textbf{脱敏处理}——
去除个人上下文，保留通用推理步骤。

在教育和客服场景中，推理链缓存的潜力尤为突出。
教育场景中，同一道题可能被数千名学生提问，
缓存后首次推理的高计算成本可以被摊薄到
几乎可忽略的水平。
客服场景中，常见问题的推理链可以被预先生成和缓存，
将平均响应延迟从数秒降至毫秒级。

\subsection{成本模型}

一个推理模型回答一个数学问题的成本构成：
\begin{itemize}[noitemsep]
  \item \textbf{标准回答}（无 thinking）：
        $\sim$500 输出 token，成本 $\sim$\$0.001。
  \item \textbf{浅层推理}（1K thinking tokens）：
        $\sim$1500 总 token，成本 $\sim$\$0.003。
  \item \textbf{深度推理}（10K thinking tokens）：
        $\sim$10500 总 token，成本 $\sim$\$0.02。
  \item \textbf{Best-of-16 + 深度推理}：
        $\sim$168K 总 token，成本 $\sim$\$0.30。
\end{itemize}

深度推理的成本是标准回答的 20 倍，
Best-of-16 + 深度推理是 300 倍。
但如果问题本身的价值足够高
（一道竞赛题的解答、一个关键 bug 的修复），
这个成本仍然极为合算。

关键在于\textbf{匹配推理投入和问题价值}——
用 \$0.30 的推理计算来回答「今天天气怎么样」是浪费，
用 \$0.001 的快速回答来解数学竞赛题是不足。
好的推理时计算系统需要\textbf{自适应}地分配资源。


\section{推理模型的前沿与未来}

\subsection{从推理到行动：Agent 范式}

2026 年的一个核心趋势是
推理模型从「思考者」进化为「行动者」。
传统推理模型只在\textbf{文本空间}中推理；
新一代模型在\textbf{行动空间}中推理——
它们可以规划行动序列、执行工具调用、
观察结果、调整计划。

这种「推理 + 行动」的范式
本质上是将推理时缩放应用于更广阔的问题空间。
不再只是「想更久」，
而是「想更久，并且在想的过程中做实验来验证想法」。

\textbf{ReAct 和 Reflexion} 是这一范式的两个重要先驱。
ReAct（Reasoning + Acting）让模型交替生成推理文本和工具调用，
形成「想→做→看→想」的循环。
Reflexion 更进一步——在任务失败后，
模型生成自我反思文本，
分析失败原因并调整策略，
然后重新尝试。
这种“从失败中学习”的能力
使 Agent 具备了类似人类的
\textbf{经验驱动的错误恢复}能力。

在代码 Agent 的实践中，
行动空间中的推理表现为多步骤的工程决策：
Agent 阅读代码 $\to$ 形成假说（“bug 可能在文件 X 的函数 Y 中”）
$\to$ 编写测试验证假说 $\to$ 运行测试
$\to$ 根据测试结果修正假说或实施修复
$\to$ 再次运行测试确认修复。
每一步都消耗 thinking tokens，
但工具执行的结果（测试输出、编译错误、运行时日志）
提供了\textbf{真实世界的反馈信号}，
比纯文本推理中的“自我对话”更可靠。

然而，行动空间中的推理也引入了新的失败模式。
\textbf{无限循环}——Agent 可能在同一类错误上
反复尝试相同的修复策略。
\textbf{工具滥用}——Agent 可能过度依赖工具调用
而非充分思考，导致大量无效的试错。
\textbf{计划偏移}——在多步骤执行中，
Agent 可能逐渐偏离原始目标，
被中间步骤的细节所吸引。
Claude Code 通过 \texttt{CLAUDE.md}
中的任务计划系统和步数上限
来缓解这些问题；
OpenAI Codex 通过 Smart Approvals
在关键决策点暂停并请求人类确认。

推理时计算在 Agent 场景中的经济学也有所不同。
传统推理模型的成本主要由 thinking tokens 决定，
而 Agent 推理的成本还包括工具执行的延迟和计算资源。
一次代码测试运行可能需要 30 秒——
这段时间 GPU 空闲，但用户在等待。
将推理 tokens 和工具执行延迟纳入统一的“计算预算”，
并在两者之间做最优分配，
是 Agent 推理经济学的核心问题。

\subsection{推理模型在 Agent 场景中的独特优势}

推理模型（R1、o3、Claude Opus 4.6）
相比非推理模型在 Agent 任务中表现出
\textbf{质的差异}，而不仅仅是量的提升。

\textbf{错误恢复能力}是最显著的差异。
非推理模型在 Agent 工作流中
遇到错误时往往“陷入循环”——
反复尝试相同的失败策略，
因为它们缺乏“回溯到更早的决策点”的能力。
推理模型通过 thinking tokens 中的自我反思
可以\textbf{诊断}失败原因（“The previous approach failed because...”），
然后\textbf{主动切换}到不同的策略。
在 SWE-bench Pro 的代码修复任务中，
推理模型的“首次尝试失败后恢复成功率”
达到 35--45\%，
而非推理模型仅为 10--15\%。

\textbf{长期规划能力}是另一个关键差异。
Agent 任务通常需要多步规划——
先分析问题 → 制定计划 → 分步执行 → 验证结果。
非推理模型倾向于“贪心执行”——
在每一步做局部最优选择，
但可能忽略了全局的一致性。
推理模型可以在 thinking tokens 中
进行“预演”（mental simulation）——
在实际执行之前\textbf{想象}多步行动的后果，
选择全局更优的路径。
这类似于象棋高手的“计算变化”——
在脑中推演多步棋后再落子。

\textbf{推理模型对 Agent 系统设计的影响}
正在改变 Agent 的架构范式。
传统 Agent 架构（如 ReAct）
依赖\textbf{外部编排}来管理推理和行动的交替——
框架代码显式地在“推理步骤”和“行动步骤”之间切换。
推理模型的 extended thinking 使得这种外部编排
变得不再必要——
模型在一个连贯的推理过程中
\textbf{自主决定}何时思考、何时行动。
这种“内化编排”的 Agent
比外部编排更灵活
（模型可以在任意时刻切换思考和行动），
但也更难调试
（推理过程对开发者不完全透明）。

Claude Code 的 Agent 架构正是这种“内化编排”的代表。
模型在 extended thinking 中同时完成
问题分析、工具选择、执行计划制定
和结果预期——
然后一次性输出包含工具调用的回复。
这种设计使得每次“思考→行动”循环
比传统 ReAct 减少了 1--2 轮 API 调用，
在延迟和成本上都有显著优势。

\subsection{效率缩放：做更多、花更少}

2026 年行业的核心议题不再是「更大的模型」，
而是\textbf{效率缩放}（efficiency scaling）——
用更少的计算达到同等甚至更好的效果。
这一范式转变可以用一个\textbf{等算力分析}（iso-compute analysis）
来量化：固定计算预算 $C$，
比较不同策略在相同成本下能达到的最优性能。

\textbf{蒸馏}是效率缩放中投入产出比最高的手段。
R1 的 7B 蒸馏模型在 MATH-500 上达到 92.8\%，
接近 70B 级别——
1/10 的参数，1/10 的推理成本。
更令人震惊的是等算力对比：
R1-Distill-Qwen-32B（AIME 72.6\%）的单次推理成本
约为 R1 全量模型（AIME 79.8\%）的 $1/20$——
如果将省下的计算用于 Best-of-8 采样，
蒸馏模型的有效准确率可以超过教师模型的单次推理。
这意味着在固定计算预算下，
\textbf{小模型 + 多次采样}可以优于\textbf{大模型 + 单次推理}。
Step 3.5 Flash 将这一策略推向极致——
仅 11B 活跃参数（196B 总参数 MoE）就达到了 AIME 97.3\%，
证明了 MoE 架构配合高质量推理蒸馏的等算力效率
可以超越参数量大 10 倍以上的稠密模型。

\textbf{混合推理模式}代表了部署效率的另一个突破。
传统方案需要部署两套模型——
一个快速模型处理简单请求，一个推理模型处理难题——
这带来了双倍的 GPU 内存和运维开销。
Qwen3~\cite{qwen2025qwen3} 的 thinking/non-thinking 统一架构
让一个模型通过\texttt{/think}和\texttt{/no\_think}标记
在两种模式间无缝切换。
在 non-thinking 模式下，模型跳过推理链直接输出答案，
延迟和成本与传统模型相当；
在 thinking 模式下，模型展开完整的推理链。
这种统一架构不仅省去了部署两套模型的成本，
还避免了路由器误判带来的质量损失——
因为路由决策可以延迟到模型\textbf{内部}完成，
而不是在外部预判问题难度。

\textbf{推理预算自适应}是效率缩放的第三个支柱。
核心洞察是：在真实的请求流量中，
80\% 以上的问题不需要深度推理。
如果对所有请求都分配相同的推理预算，
大量计算被浪费在已经能正确回答的简单问题上。
实践中的自适应策略包括：
（1）\textbf{Budget forcing}——
根据问题复杂度预分配 thinking token 上限
（简单问题 256 tokens，中等问题 4K，困难问题 32K+）；
（2）\textbf{Early exit}——
当模型在推理链中产生足够高置信度的答案时提前终止
（避免 overthinking）；
（3）\textbf{Cascading}——
先用小模型尝试回答，只在置信度低于阈值时
才升级到大模型或深度推理模式。
Google 的 Gemini 2.5 Flash 原生支持动态推理深度，
根据问题复杂度自动调整 thinking token 数量，
在不需要人工干预的前提下实现了接近最优的计算分配。

\textbf{Self-play 蒸馏}（自博弈蒸馏）
是效率缩放中的一个新兴方向。
核心思想是让模型\textbf{自己做教师}——
在困难问题上用高推理预算（64K tokens）生成正确答案，
然后用这些答案\textbf{蒸馏}回自身，
让模型学会用更短的推理链达到相同的结果。
这种“自我压缩推理”的过程可以迭代——
每轮蒸馏后，模型的推理效率提高，
再用新模型生成更好的蒸馏数据做下一轮。
初步实验表明，3--5 轮自博弈蒸馏后，
模型可以在推理链长度减少 60\% 的前提下
保持 95\% 以上的原始准确率。

\textbf{推理链压缩}（Reasoning chain compression）
从推理链本身的结构入手。
许多推理链中包含大量\textbf{冗余内容}：
重复的自我确认（“Yes, that's correct...”）、
不必要的中间步骤重述
（“So far we have established that...”）、
以及犹豫和修正的过程
（“Wait, let me reconsider... Actually, my first approach was right”）。
这些冗余在 RL 训练中是有价值的
（它们帮助模型“保持状态”不遗忘前面的推理），
但在推理时增加了不必要的计算。

两种推理链压缩方案正在被探索：
（1）\textbf{结构化推理格式}——
训练模型用更紧凑的格式输出推理
（如 JSON-like 的步骤编号结构，而非自然语言叙述），
在保持推理质量的同时将 token 数量减少 30--50\%；
（2）\textbf{隐式推理}（Latent reasoning）——
将 thinking tokens 替换为模型内部的\textbf{隐藏表示}，
不生成可读的推理文本，
而是在隐空间中进行“思考”。
这是一个更激进的方向——
如果推理过程不需要通过 token 空间“外化”，
计算效率可以提升一个数量级。
但隐式推理面临\textbf{可解释性}和\textbf{可训练性}的双重挑战：
如何验证隐空间中的“推理”是否正确？
如何为隐空间推理设计有效的 RL 奖励信号？
这些问题的答案将决定推理模型的下一代架构。

\textbf{推理链缓存}是效率缩放中尚未被充分开发的方向。
在教育和客服等场景中，大量请求涉及相似的问题，
复用已完成的推理链可以将边际成本降到接近零。
如前所述，推理链缓存面临语义匹配、缓存有效性和隐私的挑战，
但其潜在的经济价值巨大——
一道被 1000 名学生提问的数学题，
首次推理的 \$0.02 成本被摊薄为每人 \$0.00002。

综合来看，效率缩放的各种手段\textbf{可以叠加}。
一个完整的效率缩放 pipeline 可能是：
蒸馏模型 + 混合推理模式 + 难度路由 + 推理链缓存——
每一层都在上一层的基础上进一步压缩成本。
行业估计，2024 年到 2026 年间，
达到 GPT-4 等效性能的单位推理成本
下降了约 $1000\times$，
其中硬件进步贡献约 $3\times$，
软件优化（量化、batching、推测解码）贡献约 $30\times$，
效率缩放（蒸馏 + 自适应推理）贡献约 $10\times$。

\subsection{开放的问题}

推理时缩放仍有多个\textbf{未解决的问题}，
每一个都可能重新定义这一领域的方向。

\textbf{如何训练更好的 PRM？}
PRM 的质量决定了搜索策略的上限，
但获取步骤级别的训练数据极为昂贵——
标注一个推理步骤是否正确
比标注最终答案是否正确困难得多，
因为它需要标注者\textbf{理解}推理过程的每一步。
人类标注的成本使得大规模 PRM 训练几乎不可行：
OpenAI 的 PRM800K 数据集~\cite{lightman2024lets}
包含约 800K 个步骤级标注，
仅覆盖了 12K 个数学问题——
而一个鲁棒的 PRM 可能需要覆盖数百万个不同类型的推理问题。

自动化的 PRM 训练是一个活跃的研究方向。
\textbf{MCTS 回溯标注}是最有前景的路线之一：
对于一个推理步骤，
从该步骤出发多次采样到最终答案，
用最终答案的正确率来\textbf{近似}该步骤的正确性。
如果从步骤 $s_i$ 出发的 100 次采样中有 85 次得到正确答案，
则 $s_i$ 的过程奖励近似为 0.85。
这种方法不需要人类标注，但计算量巨大——
每个训练样本需要数百次的 rollout。
Math-Shepherd~\cite{wang2024mathshepherd}
使用这种方法自动构建了步骤级训练数据，
训练出的 PRM 在 GSM8K 和 MATH 上
接近甚至达到了人类标注 PRM 的水平。
另一个有趣的方向是\textbf{ORM-to-PRM 蒸馏}——
用已有的高质量 ORM（相对容易训练）
来为每个推理步骤生成伪标签，
然后训练 PRM。
这种方法的效果取决于 ORM 的质量
和推理步骤的分解粒度。

\textbf{推理时缩放有没有「上限」？}
Snell et al.\ 表明，在极难的问题上，
增加推理计算的收益会趋近于零。
但目前尚不清楚这个上限在哪里，
以及更好的推理策略能否推高这个上限。
从理论角度看，推理时缩放的上限
受到两个基本因素的约束：
（1）\textbf{模型的知识边界}——
如果模型在预训练中从未接触过某类知识，
再多的推理时计算也无法弥补知识的缺失
（你无法通过「想更久」来得出你从未学过的公式）；
（2）\textbf{搜索空间的组合爆炸}——
某些问题（如 NP-hard 问题）的搜索空间随问题规模指数增长，
有限的推理计算无法穷尽。
然而，工具增强推理可能改变这一图景——
模型可以在推理过程中通过代码执行来暴力搜索，
或通过检索来获取未见过的知识，
从而推高上限。
o3/o4-mini 在 ARC-AGI 基准上的表现
（o3-high 在 SemiPrivate 集上达到 87.5\%）
暗示工具增强可以显著扩展推理时缩放的有效范围。
一个关键的开放问题是：是否存在一个\textbf{通用的}缩放上限，
还是每类问题有不同的天花板？

\textbf{推理能力能否泛化？}
在数学和代码上训练的推理能力
能否迁移到法律推理、医学诊断或科学发现？
现有证据呈现出一幅\textbf{混合的图景}。
一方面，推理的\textbf{元策略}（meta-strategies）
确实可以迁移：
R1 和 o1 在非数学推理任务（如 GPQA 科学问答）上
也表现出显著的推理时缩放效果——
GPQA Diamond 上，o1 达到了 78.3\%，
远超非推理模型的 50--60\%。
这表明回溯、自我验证、分类讨论等推理策略
具有一定的领域通用性。

另一方面，推理的\textbf{领域知识}不能自动获得。
一个在数学上表现出色的推理模型
可能在需要深厚医学知识的诊断问题上
表现不佳——不是因为它不会推理，
而是因为它没有相关的医学知识来推理。
这暗示推理能力的泛化有两个层次：
\textbf{推理形式}的泛化（可以跨领域）
和\textbf{领域推理}的泛化（需要领域数据）。
S1~\cite{muennighoff2025s1} 的研究从另一个角度
提供了洞察——
仅用 1000 条精心挑选的推理示范就能激发
模型在多个领域的推理能力，
暗示推理能力的种子可能比预期的更小，
关键是推理训练数据的\textbf{多样性}而非数量。

\textbf{什么是推理的本质？}
当前的推理模型是否真的在「推理」，
还是只是在做更复杂的模式匹配？
这个哲学问题对实践也有深刻的影响。
如果推理模型只是学会了一套
“看起来像推理的模式”（pattern of reasoning），
那么在训练分布外的问题上，
推理能力会急剧下降——
这已经在一些实验中得到了初步验证：
当数学题的数字范围超出训练分布时，
推理模型的表现会显著退化，
即使题目的\textbf{结构}完全相同。

GSM-Symbolic~\cite{mirzadeh2024gsmsymbolic} 的实验
提供了关于这个问题的定量证据。
研究者将 GSM8K 数学题中的\textbf{数字}随机替换
（保持数学结构完全不变），
发现多个顶级 LLM 的准确率
下降了 5--15 个百分点。
更令人惊讶的是，某些“替换”甚至让题目变得\textbf{更简单}
（例如将大数字替换为小数字），
但模型的准确率仍然下降——
这暗示模型至少部分地在“记忆”训练数据中的具体数字，
而非完全从零推理。

但推理模型（R1、o1）在这类测试中的退化
\textbf{明显小于}非推理模型（GPT-4、Claude 3）——
推理模型的准确率下降约 3--8\%，
而非推理模型下降 10--20\%。
这表明 RL 训练确实赋予了模型一定的
\textbf{推理泛化能力}，
而不仅仅是更好的模式匹配。

\textbf{推理能力的“可组合性”假说}
提供了一个有前景的理论框架。
该假说认为：推理模型习得了一组
\textbf{可组合的推理原语}
（如“数学归纳”“反证法”“分类讨论”“等价变换”），
每个原语在训练数据中多次出现并被强化，
但原语之间的\textbf{组合方式}可以是新颖的。
这类似于自然语言的组合性——
你可以用有限的单词和语法规则
生成无限的新句子。
如果推理原语的“语法”足够丰富，
模型就可以在未见过的问题上
组合出有效的推理策略。

验证这个假说的一种方式是
\textbf{推理能力的迁移测试}：
如果推理原语是可组合的，
在数学领域训练的推理能力
应该能\textbf{部分迁移}到其他需要相同原语的领域
（如形式逻辑、博弈论、程序验证）。
现有证据支持这一预期——
R1 在 GPQA（科学推理）和 BQA（因果推理）上
的表现显著优于同等大小的非推理模型，
即使这些任务的领域与数学训练数据差异很大。
但迁移的程度是\textbf{有限的}——
模型在需要深厚领域知识的推理任务上
（如有机化学反应机理推导）
仍然表现不佳，
因为必要的知识原语（如化学反应规则）
在训练数据中出现得不够频繁。

\textbf{推理时缩放与“涌现”的关系}
是另一个深层次的理论问题。
传统的涌现（emergence）观察到
在训练缩放到特定阈值时，
能力突然“涌现”——
从几乎为零跳跃到较高水平。
推理时缩放是否也存在类似的\textbf{推理涌现阈值}？
即，是否存在一个推理计算量 $C^*$，
使得 $C < C^*$ 时模型完全无法解决某类问题，
$C > C^*$ 后突然能够解决？

部分证据暗示这种推理涌现确实存在。
在某些组合优化问题上，
推理链长度从 200 增加到 500 时，
成功率从 0\% 跳跃到 60\%——
而非线性增长。
一种可能的解释是：
某些问题需要\textbf{最少 $k$ 步}的推理
才能触及正确答案
（例如一个问题需要至少 3 步回溯才能找到正确路径），
不足 $k$ 步时任何部分推理都无法接近答案，
超过 $k$ 步后成功率急剧上升。
这种“推理的相变”行为
使得推理计算预算的选择变得更加微妙——
过低的预算可能完全无效
（不如直接回答），
需要达到一个\textbf{最低有效阈值}
才能发挥推理时缩放的价值。

然而，另一些证据指向了更乐观的解读。
R1-Zero 自发涌现的反证法和分类讨论
是训练数据中\textbf{未曾出现}的推理策略——
如果只是模式匹配，模型不应该「发明」新策略。
AlphaProof 在 IMO 上解出的题目
也超出了训练分布——
没有任何训练样本包含那些具体的数学竞赛题。
一种可能的综合解读是：
推理模型学会了一种“可组合的推理原语”
（composable reasoning primitives）——
每个原语本身是从训练数据中学来的模式，
但原语之间的\textbf{组合方式}是新颖的。
这类似于人类学会了加减乘除的规则（原语），
然后可以组合应用到从未见过的具体问题上。
如果这个假说成立，推理能力的泛化性
取决于模型习得的原语的\textbf{丰富度}和\textbf{可组合性}——
这为通过增加推理训练数据的多样性
来提升泛化能力指明了方向。

这些问题的答案将决定
推理时缩放能走多远，
以及它最终能否像预训练缩放那样
带来持续的、可预测的性能提升。
2026 年的初步共识是：
推理时缩放是预训练缩放的\textbf{有力互补}，
但不是\textbf{替代品}——
两者结合，加上效率缩放的工程优化，
共同定义了 LLM 能力的前沿。

\subsection{推理时缩放的“第二曲线”：从思考到行动}

回顾推理时缩放在 2024--2026 年间的发展，
可以辨识出一条清晰的\textbf{演进轨迹}。

\textbf{第一阶段}（2024 年）：\textbf{“学会思考”}。
o1 和 R1-Zero 证明了 RL 可以激发推理能力的涌现。
推理时缩放的核心价值在于“想更久 = 答案更准确”。
关注点是数学和代码等\textbf{可验证}的推理任务。

\textbf{第二阶段}（2025 年）：\textbf{“思考 + 行动”}。
o3/o4-mini 和 Claude 将工具使用集成到推理过程中。
推理时缩放的价值从“更准确的答案”
扩展到“更有效的问题解决”——
模型不仅能想，还能做。
Agent 工作流成为推理模型的主战场。

\textbf{第三阶段}（2026 年，正在发生）：\textbf{“自主推理”}。
模型不再等待人类提出问题才开始推理——
它可以\textbf{主动}识别需要解决的问题，
制定计划，持续工作，
在需要时寻求人类反馈。
这是从“回答问题的工具”到“解决问题的 Agent”的跃迁。
Claude Code 的持续 Agent 模式
和 Codex 的异步任务执行
是这一阶段的早期产品形态。

\textbf{推理时缩放与 Agent 基础设施的融合}
正在催生一个新的技术栈。
传统的推理优化（量化、batching、推测解码）
专注于\textbf{单次推理调用}的效率。
Agent 场景需要的是\textbf{多轮推理调用链}的效率——
包括推理调用之间的状态管理、
工具执行的延迟隐藏、
以及跨调用的上下文窗口管理。
这些问题催生了新的系统层抽象：
\textbf{推理会话}（reasoning session）管理
跨多次 API 调用的推理状态，
\textbf{推理缓存}不仅缓存 KV cache，
还缓存推理链的中间结论和工具调用结果，
\textbf{推理编排器}根据任务进展
动态调整推理深度和工具使用策略。

从更宏观的视角看，
推理时缩放代表了 AI 系统设计的一个\textbf{范式转变}。
传统 AI 系统的“智能”完全在\textbf{模型权重}中——
推理时只是“读取”预训练时编码的知识。
推理时缩放将一部分“智能”转移到了\textbf{推理过程}中——
模型在面对具体问题时“构建”出新的解决方案，
而不是从记忆中“检索”。
这更接近人类认知的工作方式——
我们不是从记忆中直接提取所有答案，
而是利用基础知识和推理能力\textbf{现场解决}问题。

这个范式转变也意味着
评估 AI 能力的方式需要根本性的更新。
传统的基准测试（MMLU、ARC 等）
测量的是模型的\textbf{知识存储}——
能否回忆起正确的事实。
推理时缩放时代需要的评估
更应关注模型的\textbf{推理过程}——
能否在新颖的问题上
构建出有效的解决路径。
SWE-bench、AIME、ARC-AGI
之所以成为 2025--2026 年最重要的基准，
正是因为它们测量的是后者而非前者。

推理时缩放仍然年轻——
从 o1 的发布到本文写作只有不到两年。
但它已经深刻地改变了
我们对 AI 能力边界的认知。
一个小模型“想更久”可以匹配大模型的质量，
一个推理模型“边想边做”可以完成复杂的工程任务，
一个 Agent “自主推理”可以持续地解决问题——
这些能力在两年前还不存在。
推理时缩放的下一个两年会带来什么？
如果历史有任何指导意义的话——
答案可能超出我们的想象。


%% ============================================================
%% NEW REFERENCES
%% ============================================================
% NEW REF: wang2023selfconsistency = Wang, Xuezhi, et al. "Self-Consistency Improves Chain of Thought Reasoning in Language Models." ICLR 2023.
% NEW REF: ichihara2025rbon = Ichihara, Yuki, et al. "Evaluation of Best-of-N Sampling Strategies for Language Model Alignment." TMLR 2025.
% NEW REF: park2025lemcts = Park, Sungjin, et al. "Ensembling Large Language Models with Process Reward-Guided Tree Search for Better Complex Reasoning." NAACL 2025.
% NEW REF: lu2025octotools = Lu, Pan, et al. "OctoTools: An Agentic Framework with Extensible Tools for Complex Reasoning." arXiv:2502.11271, 2025.
% NEW REF: lightman2024lets = Lightman, Hunter, et al. "Let's Verify Step by Step." ICLR 2024. arXiv:2305.20050, 2023.
% NEW REF: wang2024mathshepherd = Wang, Peiyi, et al. "Math-Shepherd: Verify and Reinforce LLMs Step-by-step without Human Annotations." ACL 2024. arXiv:2312.08935, 2024.
% NEW REF: muennighoff2025s1 = Muennighoff, Niklas, et al. "s1: Simple Test-Time Scaling." arXiv:2501.19393, 2025.
% NEW REF: qwen2025qwen3 = Qwen Team. "Qwen3 Technical Report." arXiv:2505.09388, 2025.
% NEW REF: mirzadeh2024gsmsymbolic = Mirzadeh, Iman, et al. "GSM-Symbolic: Understanding the Limitations of Mathematical Reasoning in Large Language Models." arXiv:2410.05229, 2024.
% NEW REF: demoura2008z3 = de Moura, Leonardo, and Nikolaj Bjørner. "Z3: An efficient SMT solver." TACAS 2008.
% NEW REF: hooper2024fp8 = Hooper, Coleman, et al. "KVQuant: Towards 10 Million Context Length LLM Inference with KV Cache Quantization." NeurIPS 2024.
% NEW REF: chowdhery2022palm = Chowdhery, Aakanksha, et al. "PaLM: Scaling Language Modeling with Pathways." JMLR 2023.
% NEW REF: penedo2023refinedweb = Penedo, Guilherme, et al. "The RefinedWeb Dataset for Falcon LLM." NeurIPS 2023 (Datasets Track).
